% Academic curriculum vitae -- Junho Lee
% Content source: _layouts/default.html (News excluded).
% Project summaries: projects/normal-gaussian-splats.html and
% projects/teleop-diffusion-policy.html.
% Build twice: pdflatex cv.tex

\documentclass[11pt,a4paper]{article}
\usepackage[T1]{fontenc}
\usepackage[utf8]{inputenc}
\usepackage{lmodern}
\usepackage{microtype}
\usepackage[margin=0.85in,top=0.8in,bottom=0.8in,headheight=14pt]{geometry}
\usepackage{titlesec}
\usepackage{enumitem}
\usepackage{fancyhdr}
\usepackage{lastpage}
\usepackage[hidelinks]{hyperref}

\hypersetup{
  pdftitle={Junho Lee -- Curriculum Vitae},
  pdfauthor={Junho Lee},
  pdfsubject={Academic appointments, education, publications, and research projects}
}

\setlength{\parindent}{0pt}
\setlength{\parskip}{0.25em}
\linespread{1.04}
\raggedbottom
\urlstyle{same}

% Restrained serif headings and thin rules for a traditional academic CV.
\titleformat{\section}{\large\scshape}{}{0pt}{}[{\titlerule[0.4pt]}]
\titlespacing*{\section}{0pt}{1.25em}{0.65em}
\titleformat{\subsection}{\normalsize\bfseries}{}{0pt}{}
\titlespacing*{\subsection}{0pt}{1.1em}{0.65em}

% Dates form a fixed-width column; each entry remains together on the page.
\newlength{\datewidth}
\setlength{\datewidth}{1.35in}
\newcommand{\datedentry}[3]{%
  \par\noindent
  \begin{minipage}[t]{\datewidth}
    \raggedright\small #1
  \end{minipage}\hspace{0.15in}%
  \begin{minipage}[t]{\dimexpr\linewidth-\datewidth-0.15in\relax}
    \textbf{#2}\par
    #3
  \end{minipage}\par\vspace{0.5em}}

% Bibliographic entries use hanging numbers and bold the CV author's name.
\newlist{publications}{enumerate}{1}
\setlist[publications]{label={[\arabic*]},leftmargin=2.1em,
  labelsep=0.65em,itemsep=0.85em,topsep=0pt,parsep=0pt}
\newcommand{\publication}[4]{%
  \item\begin{minipage}[t]{\linewidth}
    #1\par
    \textbf{#2}\par
    \textit{#3}\par
    {\small #4}
  \end{minipage}}
\newcommand{\resource}[2]{\href{#1}{[#2]}}

\pagestyle{fancy}
\fancyhf{}
\fancyhead[L]{\small Junho Lee}
\fancyhead[R]{\small\textsc{Curriculum Vitae}}
\renewcommand{\headrulewidth}{0.3pt}
\fancyfoot[R]{\small \thepage\ / \pageref{LastPage}}
\renewcommand{\footrulewidth}{0pt}
\fancypagestyle{firstpage}{%
  \fancyhead{}%
  \renewcommand{\headrulewidth}{0pt}}

\begin{document}
\thispagestyle{firstpage}

% Research profile and academic background.
{\LARGE\bfseries Junho Lee}\hfill{\scshape Curriculum Vitae}\par
\vspace{0.5em}
Postdoctoral Researcher\par
Autonomous Systems Lab, Institute of Computer Technology\par
TU Wien, Vienna, Austria\par
\vspace{0.3em}
{\small
\href{https://twjhlee.github.io}{twjhlee.github.io}\quad$\vert$\quad
\href{mailto:twjhlee@gmail.com}{twjhlee@gmail.com}\quad$\vert$\quad
\href{mailto:junho.lee@tuwien.ac.at}{junho.lee@tuwien.ac.at}}\par

\section{Research Profile}

My research focuses on 3D perception and affordance-based robotic manipulation.
My goal is to build reliable robotic systems that can perceive and manipulate
real-world objects in the wild.

I am a postdoctoral researcher in the
\href{https://www.tuwien.at/etit/ict/asl/team}{Autonomous Systems Lab},
\href{https://www.tuwien.at/en/etit/ict}{Institute of Computer Technology},
\href{https://www.tuwien.at/}{TU Wien}.
My principal investigator is
\href{https://www.tuwien.at/etit/ict/asl/team/dongheui-lee}{Dongheui Lee}.

\section{Academic Appointments}

\datedentry{1 Sep.\ 2026\\-- present}
{Postdoctoral Researcher}
{TU Wien, Vienna, Austria.\\
Autonomous Systems Lab, Institute of Computer Technology.}

\datedentry{1 Sep.\ 2025\\-- 31 Aug.\ 2026}
{Postdoctoral Researcher}
{Seoul National University, Seoul, Republic of Korea.}

\section{Education}

\datedentry{Sep.\ 2020\\-- Aug.\ 2025}
{Ph.D.\ in Electrical and Computer Engineering}
{Seoul National University, Seoul, Republic of Korea.\par
\textit{Dissertation:} \href{https://twjhlee.github.io/pdfs/thesis.pdf}{Toward Reliable
Vision-Based Robotic Grasping in Transparency and Clutter}.\par
\resource{https://twjhlee.github.io/pdfs/thesis.pdf}{Dissertation}}

\datedentry{Mar.\ 2017\\-- Aug.\ 2020}
{B.S.\ in Electrical and Computer Engineering}
{Seoul National University, Seoul, Republic of Korea.}

% Journal and conference publications, in reverse chronological order.
\section{Journal and Conference Publications}

\begin{publications}
\publication
  {\textbf{Junho Lee}, Sang Min Kim, Yonghyeon Lee, and Young Min Kim.}
  {INGRID: Interactive Geometry and Instance Identification for Occluded Scenes.}
  {Computers \& Graphics, 2026.}
  {\resource{https://www.sciencedirect.com/science/article/pii/S0097849326002475}{Paper}\quad
   \resource{https://twjhlee.github.io/projects/INGRID}{Project}\quad
   \resource{https://youtu.be/zDi0jirDK80}{Video}}

\publication
  {Hyeongjun Heo, Seungyeon Woo, Sang Min Kim, Junho Kim, \textbf{Junho Lee},
   Yonghyeon Lee, and Young Min Kim.}
  {AnyCamVLA: Zero-Shot Camera Adaptation for Viewpoint Robust
   Vision-Language-Action Models.}
  {IEEE/RSJ International Conference on Intelligent Robots and Systems (IROS), 2026.}
  {\resource{https://arxiv.org/abs/2603.05868}{Paper}\quad
   \resource{https://heo0224.github.io/AnyCamVLA/}{Project}}

\publication
  {\textbf{Junho Lee}, Sang Min Kim, Yonghyeon Lee, and Young Min Kim.}
  {NFL: Normal Field Learning for 6-DoF Grasping of Transparent Objects.}
  {IEEE Robotics and Automation Letters (RA-L), 2024.}
  {\resource{https://ieeexplore.ieee.org/document/10328050}{Paper}\quad
   \resource{https://twjhlee.github.io/projects/NFL}{Project}\quad
   \resource{https://youtu.be/b92CTDKeEk8}{Video}\quad
   \resource{https://github.com/twjhlee/Normal-Field-Learning}{Code}}

\publication
  {\textbf{Junho Lee}, Junhwa Hur, Inwoo Hwang, and Young Min Kim.}
  {MasKGrasp: Mask-based Grasping for Scenes with Multiple General Real-world Objects.}
  {IEEE/RSJ International Conference on Intelligent Robots and Systems (IROS), 2022.}
  {\resource{https://ieeexplore.ieee.org/document/9982130}{Paper}\quad
   \resource{https://twjhlee.github.io/projects/maskgrasp}{Project}\quad
   \resource{https://www.youtube.com/watch?v=iWvstHPu97Y}{Video}}

\publication
  {Cheol-Hui Min, Jinseok Bae, \textbf{Junho Lee}, and Young Min Kim.}
  {GATSBI: Generative Agent-centric Spatio-temporal Object Interaction.}
  {IEEE Conference on Computer Vision and Pattern Recognition (CVPR), 2021.
   Oral presentation.}
  {\resource{https://arxiv.org/abs/2104.04275}{Paper}\quad
   \resource{https://www.youtube.com/watch?v=nAf87_0T5CE&feature=youtu.be&ab_channel=3DVisionLabattheSeoulNationalUniversity}{Video}}
\end{publications}

% Preprint, ongoing projects, intellectual property, and service.
\section{Preprints}

\begin{publications}[start=6]
\publication
  {Sang Min Kim, Jinwoo Seo, Hyeongjun Heo, \textbf{Junho Lee},
   Yonghyeon Lee\textsuperscript{*}, and Young Min Kim\textsuperscript{*}.}
  {POIL: Point-based One-Shot Imitation Learning with Stable Dynamical Systems.}
  {Preprint, 2026.}
  {\textsuperscript{*}Co-corresponding authors.\quad
   \resource{https://arxiv.org/abs/2609.30404}{Paper}\quad
   \resource{https://sangminkim-99.github.io/poil/}{Project}}
\end{publications}

\section{Research Projects}

\textbf{\href{https://twjhlee.github.io/projects/normal-gaussian-splats}{Real-time
Normal Gaussian Splats}}\hfill\textit{Ongoing}\par
Real-time normal-field tracking of moving objects using Gaussian splats and SAM2 masks.

\vspace{0.65em}
\textbf{\href{https://twjhlee.github.io/projects/teleop-diffusion-policy}{Teleop
Demonstration Accumulation for Diffusion Policies}}\hfill\textit{Ongoing}\par
WebXR teleoperation of a Piper arm for demonstration collection to train diffusion policies.

\section{Patents}

\textbf{Environment-aware color measurement method, apparatus, and
computer-readable medium.}\par
Korean patent registration no.\ KR10-2596914.

\section{Honours and Awards}

\datedentry{2021}{Tutor Award}{College of Engineering, Seoul National University.}

\section{Academic Service}

\textbf{Reviewer}\par
IEEE Robotics and Automation Letters (RA-L); IEEE International Conference on
Robotics and Automation (ICRA); IEEE/RSJ International Conference on
Intelligent Robots and Systems (IROS); Pacific Graphics (PG).

\end{document}
